%
% 北見工業大学卒業論文
%
%%%%%%%%%%%%%%%%%%%%%%%%% Preamble %%%%%%%%%%%%%%%%%%%%%%%%%%%
\documentclass[twocolumn, a4paper]{ltjsarticle}
 
\usepackage{mystyle}
\usepackage{geometry}
\geometry{margin=19mm,top=25.4mm,bottom=25.4mm}

\begin{document}

%挿入するソースコードのフォント
\lstset{
  numbers=left,
  stepnumber=1,
  numberstyle=\small\rmfamily,
  basicstyle=\small\ttfamily,
  frame={tb}
}

\renewcommand{\lstlistingname}{ソースコード}

\twocolumn[
\begin{center}
{\huge{数値解析用プログラミング言語と連動したアニメーション図説フレームワークの開発に関する研究}}\\
\vspace{3mm}
{\large 西村美春}\\
\vspace{3mm}
情報デザイン・コミュニケーション工学コース　数理波動システム研究室\\
\end{center}
]

\section{背景・目的}
現在、教育資料や発表資料などのプレゼンテーション資料は、動画編集ソフトやプレゼンテーション作成ソフトを用いて作成されている。
代表的なプレゼンテーション作成ツールとして、PowerPoint\cite{Microsoft}が広く利用されており、静的な資料を作成する際には操作が容易で、視覚的にわかりやすい資料を作成できる点で優れている。
一方で、図形を視覚的操作によって作成するため、数式やグラフを用いた正確な図の作成が困難である。
また、同一または類似したアニメーションであっても、個別に作成する必要があるという問題がある。
\par
また、アニメーション作成ツールとしてManim\cite{hatomatsuManim}が挙げられる。
Manimはpythonコードを用いて数式やグラフを記述し、それらをアニメーションとして動画化できるアニメーション作成ライブラリである。
数式表現に強く、教育用途や研究発表において有用である。しかし使用にはPythonによる記述が必要であり、C言語やFortranで記述されたコードをそのまま活用することはできない。
PowerPointにおいても、数値解析で用いたコードから直接グラフやアニメーションを生成することはできず、計算結果を一度数値データとして出力し、Excel等の表計算ソフトを介して可視化する必要がある。
\par
これらの問題を解決するため、本研究では新たなアニメーション図説フレームワークを構築することを目的とする。
本研究の具体的な目的は、以下のとおりである。
  \begin{itemize}
    \item 数値解析用のプログラミング言語（Aqualis）\cite{sugisakaAqualis}と共通の記述によって、スライド資料を作成できる仕組みを実現する。
    \item 図形位置や大きさを数値によって制御可能とし、正確な図説表現を可能にする。
    \item 波形などの複雑なアニメーションを簡潔な記述で描画できる機能を提供する。
\end{itemize}

\section{使用言語}
本フレームワークにおいて使用するプログラミング言語はF\texttt{\#}である。
F\texttt{\#}は、コードが簡潔で可読性の高い関数型プログラミング言語であり、安全性が高く、複雑な処理の記述に適している。
また、基本的に変更されない定義を記述する.fsファイルとユーザ処理を記述する.fsxファイルを分離できるという特徴を持つ。
これにより、ユーザは必要最低限のコードのみを記述すればよく、記述量の削減と可読性の向上が可能となる。

\section{フレームワークの実装}
図１は、本研究で提案するアニメーション図説フレームワークにおけるアニメーション資料生成の流れを示したものである。
\inputfigure{animation_resume.pdf}{関係図}
以下では、この流れについて説明する。
\begin{lstlisting}[caption={アニメーション描画領域の生成},label={animation_fsx},breaklines=true,breakatwhitespace=false]
  html.animation outputdir "line" {sX=700; sY=440; mX=160; mY=590; backgroundColor="#bbeeff"} p (100,590) <| fun (f,p) ->
\end{lstlisting}
fsxファイルにおいて、ソースコード\ref{animation_fsx}のような記述を行うことで、fsファイルに定義されたアニメーション描画領域を生成する処理が呼び出され、描画領域が生成される。
\texttt{html.animation}は、HTML形式のアニメーション出力を生成するための関数である。
ユーザは生成されたファイルの出力先ディレクトリ、出力されるアニメーションの識別名、スライド領域の大きさ、
表示領域における基準座標、背景色、ならびにstartボタンおよびresetボタンの表示位置を指定することができる。
これらの情報をもとに、フレームワーク内部の処理が実行され、HTMLファイルおよびJavaScriptファイルが自動的に生成される。
\par
生成されたアニメーションでは、図形、テキスト、数式をコンテンツとして扱うことができる。
ユーザが各コンテンツのスタイルを指定すると、その記述に基づいてHTMLファイルが生成される。
また、コンテンツのアニメーション処理は数式によって動きを定義することができ、その記述をもとにJavaScriptファイルが生成される。
さらに、本フレームワークではアニメーション制御において、複数のコンテンツを同時に動作させることも、順番に動作させていくことも可能である。
JavaScriptファイル内で計算された座標位置がHTML側に反映されることで、アニメーションが制御されている。


\section{結果}
前述したアニメーション資料生成の流れを、ソースコード\ref{line_fsx}に示す直線のスタイルおよび動きを用いて実行すると、図2の上に示すように直線が上下に揺れるアニメーションが表示される。
\begin{lstlisting}[caption={直線のアニメーション},label={line_fsx},breaklines=true,breakatwhitespace=false]
let line0 = f.animationLine Style[stroke.width 3.0; stroke.color "#ff4c4c"]
f.seq {FrameTime=30; FrameNumber=1000} <| fun _ ->
  line0.P {Start = {X = fun _ -> Dbl_e 200
                    Y = fun t -> 240+100*asm.sin(-2.0*Math.PI*t/100)}
           End = {X = fun _ -> Dbl_e 400
                  Y = fun t -> 240+100*asm.sin(2.0*Math.PI*t/100)}}
\end{lstlisting}
ここでは、正弦関数を用いて直線の始点および終点のY座標を時間に応じて変化させることで、滑らかな上下運動を実現している。
\par
また、図2の下に示すアニメーションは繰り返し構文を用いて同一の図形を複数生成し、それらを同時に動作させた例である。
このように、本フレームワークでは簡潔な記述によって、複数の図形を用いた動的なアニメーションを容易に作成することができる。
\inputfigure{line.pdf}{図形のアニメーション}

\section{まとめ}
本研究では、アニメーション図説フレームワークを開発した。
本フレームワークにより、数値解析用プログラミング言語（Aqualis）と共通のコードを用いて、図形や数式を含むアニメーションスライドを作成することが可能となった。
これにより、解析結果の視覚的な説明を効率的に行う手法を示した。

\begin{thebibliography}{99}
    % main.bblからコピーする
    \bibitem{Microsoft} Microsoft, ``Microsoft powerpoint.'' https://www.microsoft.com/ja-jp/microsoft-365/powerpoint.
    \bibitem{hatomatsuManim} はとまつ, ``理系のためのmanim入門{*}随時更新{*}.'' https://note.com/hatomatsu/n/ncc87dc37a976.
    \bibitem{sugisakaAqualis} 杉坂純一郎, ``{Sugisaka/Aqalis}.'' https://github.com/Sugisaka/Aqualis.
    %\bibitem{label1} 著者, "タイトル", 誌名, vol. ○○, pp. ○○--○○, year.
\end{thebibliography}

\end{document}
